\begin{pastexam}[label={px:gk3-112-ds-02}]{高考三級資訊處理 112 年 第二題}
（申論；25 分）

某一物流公司有下圖所示的8個地點要運送，每條方向性連線及其數字代表兩個地點的運送順序及運送成本。
\begin{center}
\begin{tikzpicture}[x=1.15cm,y=1.15cm,v/.style={circle,draw,minimum size=7mm,inner sep=0pt},
  >={Triangle[length=2.2mm,width=2mm]},w/.style={font=\small,inner sep=1.5pt}]
\node[v] (1) at (0,0) {1};
\node[v] (3) at (1.3,1.2) {3};
\node[v] (2) at (1.3,-1.2) {2};
\node[v] (4) at (2.6,0) {4};
\node[v] (5) at (3.9,1.2) {5};
\node[v] (6) at (3.9,-1.2) {6};
\node[v] (8) at (5.4,0.75) {8};
\node[v] (7) at (5.4,-0.55) {7};
\draw[->] (1)--node[w,above left]{4}(3);
\draw[->] (1)--node[w,below left]{4}(2);
\draw[->] (3)--node[w,left,pos=0.45]{3}(2);
\draw[->] (3)--node[w,left,pos=0.55]{2}(4);
\draw[->] (4)--node[w,right,pos=0.6]{3}(2);
\draw[->] (4)--node[w,above left,pos=0.6]{2}(5);
\draw[->] (4)--node[w,left,pos=0.6]{2}(6);
\draw[->] (2)--node[w,below]{5}(6);
\draw[->] (6)--node[w,right,pos=0.55]{2}(5);
\draw[->] (6)--node[w,below right]{3}(7);
\draw[->] (5)--node[w,below,pos=0.5]{3}(8);
\draw[->] (5)--node[w,left,pos=0.65]{2}(7);
\draw[->] (7)--node[w,right]{4}(8);
\end{tikzpicture}
\end{center}
\begin{itemize}[leftmargin=2.2em]
\item[㈠] 試使用拓樸排序法，找出此8個地點的運送順序以及總共運送成本。（15分）
\item[㈡] 若將上圖的地點2與地點4之間以及地點6與地點7之間的連線方向顛倒，則運用拓樸排序法後，此8個地點的運送順序以及總共運送成本為何？（10分）
\end{itemize}
\end{pastexam}

\begin{pastsoln}
答案：㈠ 運送順序 $1, 3, 4, 2, 6, 5, 7, 8$，總成本 $22$；㈡ 顛倒後圖上有環 $5 \to 7 \to 6 \to 5$，拓樸排序只能排出
$1, 3, 2, 4$ 就卡住，不存在 $8$ 個地點的運送順序。

\textbf{解讀}。 箭頭 $u \to v$ 表示 $u$ 要排在 $v$ 前面送。 「總共運送成本」理解成依運送順序從一個地點走到下一個
地點、所走的那些連線的成本總和。 ㈠ 的拓樸順序剛好是唯一的，而且相鄰兩點之間都有連線，所以這個解讀可行。 拓樸
排序用 Kahn 演算法（資結書第~22.1~節，虛擬碼~22.1）：每次取出一個入分支度為 $0$ 的點，把它的出邊刪掉。

\textbf{㈠}。 圖有 $13$ 條邊：$1 \to 3$（4）、$1 \to 2$（4）、$3 \to 2$（3）、$3 \to 4$（2）、$4 \to 2$（3）、$4 \to 5$（2）、
$4 \to 6$（2）、$2 \to 6$（5）、$6 \to 5$（2）、$6 \to 7$（3）、$5 \to 8$（3）、$5 \to 7$（2）、$7 \to 8$（4）。 起初各點的入分支度：
\begin{center}
\begin{tabular}{c*{8}{c}}
\toprule
點 & 1 & 2 & 3 & 4 & 5 & 6 & 7 & 8 \\
\midrule
入分支度 & 0 & 3 & 1 & 1 & 2 & 2 & 2 & 2 \\
\bottomrule
\end{tabular}
\end{center}
\begin{center}
\small
\begin{tabular}{clll}
\toprule
步驟 & 取出 & 刪邊後入分支度變 $0$ 的點 & 候選（入分支度 $0$） \\
\midrule
1 & 1 & 3（2 的入分支度 $3 \to 2$） & $\{3\}$ \\
2 & 3 & 4（2：$2 \to 1$） & $\{4\}$ \\
3 & 4 & 2（2：$1 \to 0$；5：$2 \to 1$；6：$2 \to 1$） & $\{2\}$ \\
4 & 2 & 6 & $\{6\}$ \\
5 & 6 & 5（7：$2 \to 1$） & $\{5\}$ \\
6 & 5 & 7（8：$2 \to 1$） & $\{7\}$ \\
7 & 7 & 8 & $\{8\}$ \\
8 & 8 & & $\varnothing$ \\
\bottomrule
\end{tabular}
\end{center}
每一步候選都只有一個點，所以拓樸順序是唯一的：$1, 3, 4, 2, 6, 5, 7, 8$。 （拓樸順序唯一，等價於順序中相鄰
兩點之間都有邊，資結書定理~22.5。）相鄰兩點之間的連線與成本：
\[
1 \xrightarrow{4} 3 \xrightarrow{2} 4 \xrightarrow{3} 2 \xrightarrow{5} 6 \xrightarrow{2} 5 \xrightarrow{2} 7 \xrightarrow{4} 8,
\qquad 4 + 2 + 3 + 5 + 2 + 2 + 4 = 22.
\]
總共運送成本 $22$。 （若把「總成本」讀成圖上所有連線的成本和，則是 $39$，但那和運送順序無關，不合題意。）

\textbf{㈡}。 顛倒後 $4 \to 2$ 變成 $2 \to 4$（成本 $3$），$6 \to 7$ 變成 $7 \to 6$（成本 $3$）。 這時
$5 \to 7$、$7 \to 6$、$6 \to 5$ 形成一個有向環。 有向環上的點互相要求排在對方前面，不可能都排進拓樸順序
（資結書定理~22.2：有拓樸順序 $\iff$ 沒有有向環）。 實際跑 Kahn 演算法：起初入分支度 $1{:}\,0$、$2{:}\,2$、
$3{:}\,1$、$4{:}\,2$、$5{:}\,2$、$6{:}\,3$、$7{:}\,1$、$8{:}\,2$。
\begin{center}
\small
\begin{tabular}{clll}
\toprule
步驟 & 取出 & 刪邊後的變化 & 候選 \\
\midrule
1 & 1 & 3：$1 \to 0$；2：$2 \to 1$ & $\{3\}$ \\
2 & 3 & 2：$1 \to 0$；4：$2 \to 1$ & $\{2\}$ \\
3 & 2 & 4：$1 \to 0$；6：$3 \to 2$ & $\{4\}$ \\
4 & 4 & 5：$2 \to 1$；6：$2 \to 1$ & $\varnothing$ \\
\bottomrule
\end{tabular}
\end{center}
只排出 $1, 3, 2, 4$，剩下的 $5, 6, 7, 8$ 入分支度都不是 $0$（$5, 6, 7$ 各是 $1$，來自環上的邊；$8$ 是 $2$），
演算法停在這裡，偵測到圖有環。 所以顛倒之後\textbf{沒有}能把 $8$ 個地點都排進去的運送順序，也就沒有總運送成本；
能確定的只有前四個地點 $1 \to 3 \to 2 \to 4$，成本 $4 + 3 + 3 = 10$。 要能運送，必須先把環 $5 \to 7 \to 6 \to 5$
上的某一條連線拿掉或改向。
\end{pastsoln}
